%
% Front matter
%

\begin{frontmatter}

\title{
  T2K 6.5700$\times10^{20}$-POT 3-Flavour Muon-Neutrino Disappearance Analysis
}

\vspace*{0.7cm}

\author[RAL,LIV]  {Costas Andreopoulos}
\author[ETHZ]     {Fatih Bay}
\author[RAL,OXF]  {Thomas Dealtry}
\author[RAL,WAR]  {Steve Dennis}
\author[IFIC]     {Lorena Escudero S\'anchez}
\author[LANCS]    {\\Nick Grant}
\author[RAL,OXF]  {and Alfons Weber}

\address[RAL]   {Science \& Technology Facilities Council (STFC), Rutherford Appleton Laboratory,\\ Harwell Oxford Campus, OX11 0QX, UK}
\address[LIV]   {University of Liverpool, Physics Department,\\ Liverpool, L69 7ZE, UK} 
\address[ETHZ]  {ETH Zurich, Institute for Particle Physics,\\ 8092 Zurich, Switzerland}
\address[OXF]   {University of Oxford, Subdepartment of Particle Physics,\\ Oxford, OX1 3RH, UK}
\address[WAR]   {University of Warwick, Physics Department,\\ Coventry, CV4 7AL, UK}
\address[IFIC]  {CSIC and University of Valencia, Instituto de Fisica Corpuscular, \\ E-46071 Valencia, Spain}
\address[LANCS] {Lancaster University, Physics Department,\\ Lancaster, LA1 4YB, UK}

%
% Abstract
%

\begin{abstract}

We report the results of the
3-flavour $\nu_{\mu}$-disappearance analysis, performed by the {\em VALOR} group,
on the combined Run 1+2+3+4 dataset.
This analysis is an update of the work presented in Refs. 
\cite{OurRun1DisappAnalysis},
\cite{OurRun12DisappAnalysis}, 
\cite{OurRun12DisappAnalysis3Flavour},
\cite{OurRun123DisappAnalysis3Flavour} and
\cite{OurRun1234aDisappAnalysis3Flavour}
with improved methods and an enlarged dataset.
The Run 1+2+3+4 dataset corresponds to an integrated J-PARC neutrino beam
exposure of 6.5700$\times$10$^{20}$ POT.

This analysis predicts \NskNoOscExpectation $\pm$ \NskNoOscExpectationError (syst)
1 $\mu$-like ring events in SuperK in the absence of any oscillation,
but only \NskObserved were observed. The observed deficit has a strong energy dependence;
the ratio of observed to expected, under the no-oscillation hypothesis, is
$\sim$\NskNoOscRatioLow\% $<$ 0.5 GeV, 
$\sim$\NskNoOscRatioMid\% between 0.5 and 1 GeV and 
$\sim$\NskNoOscRatioHig\% $>$ 1 GeV.

A $\nu_{\mu}$-disappearance analysis was performed 
in a framework of 3-flavour oscillations including matter effects in constant-density matter.
The observed reconstructed energy spectrum of 1 $\mu$-like ring events was fitted, 
and separate fits were made for the normal and the inverted mass hierarchies. In these fits, 
$\sin^{2}\theta_{23}$ and either $|\Delta m^{2}_{32}|$ (normal hierarchy) or $|\Delta m^{2}_{31}|$ (inverted hierarchy) 
were allowed to float.
The oscillation parameters $\sin^{2}\theta_{12}$, $\sin^{2}\theta_{13}$, $\Delta m^{2}_{21}$ and $\delta_{CP}$
and all 41 systematic parameters considered in this analysis were also allowed to float in the fit.

The 3-flavour $\nu_{\mu}$-disappearance fit of the combined Run 1+2+3+4 dataset,
assuming the normal mass hierarchy,
gives
\resultNormalDmsq and \resultNormalMixing
as the best-fit values (p-value = \resultNormalPvalue).
Assuming the inverted mass hierarchy, the fit yields
\resultInvertedDmsq and \resultInvertedMixing
as the best-fit values (p-value = \resultInvertedPvalue).

Confidence regions in the parameter spaces $\sin^{2}\theta_{23}$ - $|\Delta m^{2}_{32}|$ (normal hierarchy) and
$\sin^{2}\theta_{23}$ - $|\Delta m^{2}_{31}|$ (inverted hierarchy) 
are shown using both the constant-$\Delta\chi^2$ method and the method of Feldman and Cousins.
We use a new procedure that is an extension of the Feldman-Cousins method to obtain confidence intervals in 
$\sin^{2}\theta_{23}$ and either $|\Delta m^{2}_{32}|$ (normal hierarchy) or $|\Delta m^{2}_{31}|$ (inverted hierarchy) 
separately. The 68\% (90\%) confidence intervals in these parameters can be summarized as
% This is 1D FC
\resultNormalMixingOneSigmaFCLo\xspace(\resultNormalMixingTwoSigmaFCLo)
$< \sin^{2}\theta_{23} <$
\resultNormalMixingOneSigmaFCHi\xspace(\resultNormalMixingTwoSigmaFCHi)
and
\resultNormalDmsqOneSigmaFCLo$\times$10$^{-3}$ (\resultNormalDmsqTwoSigmaFCLo$\times$10$^{-3}$) eV$^{2}$/c$^{4}$
$< |\Delta m^{2}_{32}| < $
\resultNormalDmsqOneSigmaFCHi$\times$10$^{-3}$ (\resultNormalDmsqTwoSigmaFCHi$\times$10$^{-3}$) eV$^{2}$/c$^{4}$
for the normal mass hierarchy,
and
\resultInvertedMixingOneSigmaFCLo\xspace(\resultInvertedMixingTwoSigmaFCLo)
$< \sin^{2}\theta_{23} <$
\resultInvertedMixingOneSigmaFCHi\xspace(\resultInvertedMixingTwoSigmaFCHi)
and
\resultInvertedDmsqOneSigmaFCLo$\times$10$^{-3}$ (\resultInvertedDmsqTwoSigmaFCLo$\times$10$^{-3}$) eV$^{2}$/c$^{4}$
$< |\Delta m^{2}_{31}| < $
\resultInvertedDmsqOneSigmaFCHi$\times$10$^{-3}$ (\resultInvertedDmsqTwoSigmaFCHi$\times$10$^{-3}$) eV$^{2}$/c$^{4}$
for the inverted mass hierarchy.
% % This is const dchisq
% \resultNormalMixingOneSigmaLo\xspace(\resultNormalMixingTwoSigmaLo)
% $< \sin^{2}\theta_{23} <$
% \resultNormalMixingOneSigmaHi\xspace(\resultNormalMixingTwoSigmaHi)
% and
% \resultNormalDmsqOneSigmaLo$\times$10$^{-3}$ (\resultNormalDmsqTwoSigmaLo$\times$10$^{-3}$) eV$^{2}$/c$^{4}$
% $< |\Delta m^{2}_{32}| < $
% \resultNormalDmsqOneSigmaHi$\times$10$^{-3}$ (\resultNormalDmsqTwoSigmaHi$\times$10$^{-3}$) eV$^{2}$/c$^{4}$
% for the normal mass hierarchy,
% and
% \resultInvertedMixingOneSigmaLo\xspace(\resultInvertedMixingTwoSigmaLo)
% $< \sin^{2}\theta_{23} <$
% \resultInvertedMixingOneSigmaHi\xspace(\resultInvertedMixingTwoSigmaHi)
% and
% \resultInvertedDmsqOneSigmaLo$\times$10$^{-3}$ (\resultInvertedDmsqTwoSigmaLo$\times$10$^{-3}$) eV$^{2}$/c$^{4}$
% $< |\Delta m^{2}_{31}| < $
% \resultInvertedDmsqOneSigmaHi$\times$10$^{-3}$ (\resultInvertedDmsqTwoSigmaHi$\times$10$^{-3}$) eV$^{2}$/c$^{4}$
% for the inverted mass hierarchy.


\end{abstract}

%
% Keywords
%

%\begin{keyword}
%\end{keyword}

\end{frontmatter}

\clearpage

\tableofcontents

\clearpage

%\listoffigures
%\listoftables


